% chktex-file 8
% chktex-file 12
% chktex-file 18
\beginsong{Kometa}[by={Jaromír Nohavica}]

\beginverse{}
Sp\ch{Ami}{}{}{}atřil jsem kometu, oblohou letěla,
chtěl jsem jí zazpívat, ona mi zmizela,
zm\ch{Dmi}{}{}{}izela jako laň \ch{G7}{}{}{}u lesa v remízku,
v \ch{C}{}{}{}očích mi zbylo jen p\ch{E7}{}{}{}ár žlutých penízků.
\endverse{}

\beginverse{}
Penízky ukryl jsem do hlíny pod dubem,
až příště přiletí, my už tu nebudem,
my už tu nebudem, ach, pýcho marnivá,
spatřil jsem kometu, chtěl jsem jí zazpívat.
\endverse{}

\beginchorus{}
\ch{Ami}{}{}{}O vodě, o trávě, \ch{Dmi}{}{}{}o lese,
\ch{G7}{}{}{}o smrti, se kterou smířit n\ch{C}{}{}{}ejde se, \ch{E}{}{}{}
\ch{Ami}{}{}{}o lásce, o zradě, \ch{Dmi}{}{}{}o světě
\ch{E}{}{}{}a o všech lidech, co kd\ch{E7}{}{}{}y žili na téhle pl\ch{Ami}{}{}{}anetě.
( \ch{Dmi}{}{}{} \ch{E}{}{}{} )
\endchorus{}

\beginverse{}
Na hvězdném nádraží cinkají vagóny,
pan Kepler rozepsal nebeské zákony,
hledal, až nalezl v hvězdářských triedrech
tajemství, která teď neseme na bedrech.
\endverse{}

\beginverse{}
Velká a odvěká tajemství přírody,
že jenom z člověka člověk se narodí,
že kořen s větvemi ve strom se spojuje
a krev našich nadějí vesmírem putuje.
\endverse{}

\beginchorus{}
Na na na \ldots{}
\endchorus{}

\beginverse{}
Spatřil jsem kometu, byla jak reliéf
zpod rukou umělce, který už nežije,
šplhal jsem do nebe, chtěl jsem ji osahat,
marnost mne vysvlékla celého donaha.
\endverse{}

\beginverse{}
Jak socha Davida z bílého mramoru
stál jsem a hleděl jsem, hleděl jsem nahoru,
až příště přiletí, ach, pýcho marnivá,
já už tu nebudu, ale jiný jí zazpívá.
\endverse{}

\beginchorus{}
O vodě, o trávě, o lese,
o smrti, se kterou smířit nejde se,
o lásce, o zradě, o světě,
bude to písnička o nás a kometě \ldots{}
\endchorus{}

\endsong{}
